% Generated by LogicInteractionExample.lean from LeanFM.Logic.run; do not edit.
\medskip\noindent\textbf{A never-claim that fails (\texttt{check\_ctl})}\par\nopagebreak
\begin{lstlisting}[language={},breaklines=true]
{"task": "mcp_logic_consult",
 "formula": {"never": {"after": "LogicInputRejected"}}}
\end{lstlisting}
\noindent returns
\begin{lstlisting}[language={},breaklines=true]
{"verdict": "disproved",
 "tool": "check_ctl",
 "task": "mcp_logic_consult",
 "reading":
 "on every path, never LogicInputRejected was just delivered",
 "reachableNodes": 4,
 "mode": "never",
 "formula": {"never": {"after": "LogicInputRejected"}},
 "evidence":
 {"path":
  [{"state": "drafting", "message": null},
   {"state": "interpreting", "message": "LogicConsultRequested"},
   {"state": "rejected", "message": "LogicInputRejected"}],
  "kind": "counterexample"},
 "deadlocks": []}
\end{lstlisting}

\medskip\noindent\textbf{A requirement over hidden state (\texttt{classify\_requirement})}\par\nopagebreak
\begin{lstlisting}[language={},breaklines=true]
{"task": "mcp_logic_consult", "predicates": ["the cache is consistent"]}
\end{lstlisting}
\noindent returns
\begin{lstlisting}[language={},breaklines=true]
{"verdict": "classified",
 "tool": "classify_requirement",
 "reading":
 "some atoms are implementation-only; restate them over proposed events before the requirements gate",
 "atoms":
 [{"restated":
   {"AG":
    {"implies":
     [{"after": "CacheConsistentObserved"},
      {"field":
       {"value": "true",
        "name": "holds",
        "message": "CacheConsistentObserved"}}]}},
   "reason":
   "no declared message or state reveals this; no event trace can test it",
   "proposedEvent":
   {"to": "an observer that can check it",
    "message": "CacheConsistentObserved",
    "from": "the actor that owns this state",
    "fields": ["subject", "holds", "observed_at"]},
   "class": "implementation-only",
   "atom": "the cache is consistent",
   "advice":
   "move it to Implementation.lean as an internal invariant, or restate it over an event that exposes it"}]}
\end{lstlisting}

\medskip\noindent\textbf{Designing a sequence diagram by writing events (\texttt{check\_trace})}\par\nopagebreak
\begin{lstlisting}[language={},breaklines=true]
{"events":
 [{"task": "mcp_logic_consult",
   "src": "Agent",
   "message": "LogicConsultRequested",
   "id": "e1",
   "dst": "LeanFMServer"},
  {"task": "mcp_logic_consult",
   "src": "LeanFMServer",
   "prevIds": ["e1"],
   "message": "LogicVerdict",
   "id": "e2",
   "dst": "Agent",
   "causeIds": ["e1"]}]}
\end{lstlisting}
\noindent returns
\begin{lstlisting}[language={},breaklines=true]
{"verdict": "admissible",
 "tool": "check_trace",
 "reading": "2 events over 1 task(s); 0 concurrent pair(s)",
 "issues": [],
 "concurrentCount": 0,
 "concurrent": [],
 "acceptance":
 [{"task": "mcp_logic_consult",
   "states": ["answered"],
   "reachedTerminal": true,
   "accepted": true}]}
\end{lstlisting}
\noindent with the generated Mermaid sequence diagram
\begin{lstlisting}[language={},breaklines=true]
sequenceDiagram
    participant Agent
    participant LeanFMServer
    Agent->>LeanFMServer: e1 LogicConsultRequested
    LeanFMServer->>Agent: e2 LogicVerdict
\end{lstlisting}

